\chapter{Forecast Evaluation}
\label{ch:evaluation}

\begin{application}[Why This Chapter Is Non-Negotiable]
This chapter teaches the evaluation methodology used across all project directions.
Every volatility forecast you produce must be evaluated with $\QLIKE$ (not MSE),
compared with the Diebold--Mariano test, and placed in a Model Confidence Set.
If you use cross-validation, it must be purged.
If you report Sharpe ratios, they must be deflated.
These are not optional extras; they are the minimum standard for credible work.
Skip this chapter and your results mean nothing.
\end{application}

A good forecast is useless if you cannot prove it is good.
This chapter gives you the statistical machinery to distinguish genuine forecasting ability from noise, luck, and overfitting.
We cover the right loss function ($\QLIKE$), the right comparison test (Diebold--Mariano), the right multi-model framework (Model Confidence Set), the right cross-validation (purged K-fold with embargo), and the right Sharpe ratio adjustment (Deflated Sharpe Ratio).


% ══════════════════════════════════════════════════════════════
\section{Why Evaluation Methodology Matters}
\label{sec:eval-why}
% ══════════════════════════════════════════════════════════════

Before diving into specific tools, consider two scenarios that illustrate why evaluation methodology \emph{is} the credibility of your results.

\textbf{Scenario 1.}
You build a LightGBM volatility forecast that achieves 5\% lower $\QLIKE$ than the HAR benchmark (Chapter~\ref{ch:har}).
Your manager asks: ``Is that improvement statistically significant, or would it vanish on a different sample?''
Without a Diebold--Mariano test, you cannot answer.

\textbf{Scenario 2.}
You try 30 different feature sets, pick the best one, and report a backtest Sharpe ratio of 1.5.
A colleague asks: ``What's the probability that at least one of 30 random strategies would have produced a Sharpe that high?''
Without the Deflated Sharpe Ratio, you cannot answer.

\begin{warning}[The Two Failure Modes]
Evaluation errors come in two flavors:
\begin{enumerate}[nosep]
  \item \textbf{Declaring a winner that isn't one.} A 5\% improvement in $\QLIKE$ that is not statistically significant is noise, not signal.
  \item \textbf{Reporting a Sharpe ratio inflated by multiple testing.} A Sharpe of 1.5 from 30 experiments may be pure luck \citep{Bailey2014DSR}.
\end{enumerate}
\end{warning}

The evaluation framework is not a final step you tack on after research.
It is the infrastructure you build \emph{first}, so that every experiment you run produces honest, comparable numbers from day one.


% ══════════════════════════════════════════════════════════════
\section{MSE and Its Limitations for Volatility}
\label{sec:eval-mse}
% ══════════════════════════════════════════════════════════════

We start with the loss function you already know, then explain why it is not enough for volatility forecasting.

\begin{prereq}[Loss Functions]
A \textbf{loss function} $L(\sigma^2_t, h_t)$ measures how far a forecast $h_t$ is from the truth $\sigma^2_t$.
Lower loss means a better forecast.
You have used mean squared error (MSE) in regression problems; it is the default in most ML pipelines.
\end{prereq}

\subsection{The MSE Formula}

The mean squared error between a sequence of forecasts $\{h_t\}$ and realized values $\{\sigma^2_t\}$ is:
\begin{equation}
\label{eq:mse}
\text{MSE} = \frac{1}{T}\sum_{t=1}^{T} \bigl(\sigma^2_t - h_t\bigr)^2
\end{equation}
where:
\begin{itemize}[nosep]
  \item $\sigma^2_t$ is the true (unobservable) variance on day $t$,
  \item $h_t$ is the forecast variance for day $t$, produced before day $t$,
  \item $T$ is the number of forecast evaluation days.
\end{itemize}

\begin{intuition}[In Plain English]
MSE measures the average squared gap between what you predicted and what actually happened.
Squaring means big misses count disproportionately: a forecast that is off by 2 units contributes 4 to the loss, while being off by 4 units contributes 16.
It treats over-prediction and under-prediction identically.
\end{intuition}

\subsection{The Proxy Problem}

We never observe true variance $\sigma^2_t$.
Instead, we use a proxy: realized variance $\RV_t$ (Chapter~\ref{ch:rv}).
This proxy is noisy because $\RV_t = \sigma^2_t + \eta_t$, where $\eta_t$ is measurement error from finite sampling, microstructure noise (Chapter~\ref{ch:noise}), and jumps (Chapter~\ref{ch:jumps}).

The good news: MSE produces correct model \emph{rankings} even when using a noisy proxy, as long as the noise is independent of the forecast \citep{Patton2011}.
If model A has lower MSE than model B when evaluated against $\RV_t$, the same ranking holds against true $\sigma^2_t$.
This property is called \textbf{robustness to noise in the proxy}.

\subsection{Why MSE Is Still Not Enough}

MSE has a deeper problem: it is symmetric and heavily penalizes extreme values.

\begin{intuition}[MSE Penalizes Outliers Disproportionately]
Suppose your forecast is $h_t = 1.0$ (in annualized variance units) for every day.
On 249 normal days, true variance is 1.0 and MSE contribution is 0.
On one crisis day, true variance spikes to 10.0, and MSE contribution is $(10 - 1)^2 = 81$.
That single day dominates the entire loss.
This means MSE-optimal forecasts chase outliers: they overweight extreme days at the expense of forecasting accuracy during normal times, which is the opposite of what most applications need.
\end{intuition}

\begin{keyidea}[MSE Is Necessary but Not Sufficient]
MSE is proxy-robust, which is valuable.
But its sensitivity to extreme realized variance values makes it a poor primary metric for volatility.
Report it as a secondary check alongside $\QLIKE$.
\end{keyidea}

\begin{projectconnection}[Why This Matters]
In your HAR vs.\ ML comparison, MSE will be dominated by a handful of crisis days (e.g., COVID March 2020, VIX spikes).
A model that slightly better predicts those extremes will look dramatically better under MSE, even if it performs worse on the 95\% of normal days that matter for daily risk management.
Always report MSE as a secondary metric, but never let it drive model selection.
\end{projectconnection}


% ══════════════════════════════════════════════════════════════
\section{QLIKE: The Preferred Loss Function}
\label{sec:eval-qlike}
% ══════════════════════════════════════════════════════════════

Having seen why MSE over-penalizes extreme days, we now introduce the loss function that the volatility forecasting literature has converged on as the primary metric.

\subsection{Intuition}

$\QLIKE$ (quasi-likelihood loss) comes from the negative log-likelihood of a Gaussian distribution with variance $h_t$.
Think of it this way: if returns were exactly normal with variance $h_t$, the best forecast would minimize $\QLIKE$.
Even when returns are not normal (and they never are), $\QLIKE$ retains two critical properties that MSE shares one of and lacks the other.

\subsection{The QLIKE Formula}

\begin{equation}
\label{eq:qlike}
\QLIKE = \frac{1}{T}\sum_{t=1}^{T} \left(\ln h_t + \frac{\sigma^2_t}{h_t}\right)
\end{equation}
where:
\begin{itemize}[nosep]
  \item $h_t$ is the forecast variance for day $t$,
  \item $\sigma^2_t$ is the true variance on day $t$ (in practice, $\RV_t$),
  \item $\ln h_t$ penalizes forecasts that are too large (over-prediction),
  \item $\sigma^2_t / h_t$ penalizes forecasts that are too small (under-prediction),
  \item $T$ is the number of evaluation days.
\end{itemize}

\begin{intuition}[In Plain English]
QLIKE has two parts pulling in opposite directions.
The $\ln h_t$ term punishes you for forecasting too high (wasting capital on unnecessary hedges), while the $\sigma^2_t / h_t$ term punishes you for forecasting too low (holding unrecognized risk).
Critically, the under-prediction penalty ($\sigma^2_t / h_t$) explodes as $h_t \to 0$, so QLIKE is far harsher on dangerous under-estimates than on conservative over-estimates.
This asymmetry matches real-world priorities: underestimating volatility gets you fired; overestimating it merely costs some opportunity.
\end{intuition}

\begin{intuition}[Why QLIKE Is Less Sensitive to Outliers]
When true variance spikes to $\sigma^2_t = 10$ and your forecast is $h_t = 1$, the QLIKE contribution is $\ln(1) + 10/1 = 10$.
Under MSE, the same day contributes $(10 - 1)^2 = 81$.
QLIKE penalizes the error linearly (through the ratio $\sigma^2_t / h_t$) rather than quadratically.
Extreme days still matter, but they do not dominate.
\end{intuition}

\begin{keyresult}[Patton (2011): QLIKE and MSE Are the Only Robust Losses]
\citet{Patton2011} proves that QLIKE and MSE are the \emph{only} two members of the standard loss function family that produce correct model rankings even when the volatility proxy is noisy.
Other common losses (MAE, HMSE, heteroskedasticity-adjusted MSE) can reverse the true ranking when evaluated against $\RV_t$ instead of $\sigma^2_t$.
Of the two robust losses, QLIKE is less sensitive to extreme $\RV$ days and is therefore preferred as the primary evaluation metric.
\end{keyresult}

\subsection{Worked Example: MSE vs.\ QLIKE Rankings}

\begin{workedexample}{MSE vs.\ QLIKE Can Disagree}
Consider five days of forecasts from two models, evaluated against realized variance:

\medskip
\begin{center}
\begin{tabular}{cccc}
\toprule
Day & $\RV_t$ & Model A ($h^A_t$) & Model B ($h^B_t$) \\
\midrule
1 & 1.2 & 1.0 & 1.3 \\
2 & 0.9 & 1.0 & 0.8 \\
3 & 1.1 & 1.0 & 1.0 \\
4 & 1.0 & 1.0 & 1.1 \\
5 & 8.0 & 1.0 & 3.0 \\
\bottomrule
\end{tabular}
\end{center}
\medskip

Model A is a constant forecast ($h_t = 1.0$).
Model B adapts but is still far from the spike on day 5.

\textbf{MSE computation:}
\begin{align*}
\text{MSE}_A &= \frac{1}{5}\bigl[(1.2-1)^2 + (0.9-1)^2 + (1.1-1)^2 + (1-1)^2 + (8-1)^2\bigr] \\
  &= \frac{1}{5}(0.04 + 0.01 + 0.01 + 0 + 49) = 9.812 \\[4pt]
\text{MSE}_B &= \frac{1}{5}\bigl[(1.2-1.3)^2 + (0.9-0.8)^2 + (1.1-1)^2 + (1-1.1)^2 + (8-3)^2\bigr] \\
  &= \frac{1}{5}(0.01 + 0.01 + 0.01 + 0.01 + 25) = 5.008
\end{align*}

MSE ranks Model B first (5.008 < 9.812), driven almost entirely by day 5.

\textbf{QLIKE computation:}
\begin{align*}
\QLIKE_A &= \frac{1}{5}\left[\left(\ln 1 + \frac{1.2}{1}\right) + \left(\ln 1 + \frac{0.9}{1}\right) + \left(\ln 1 + \frac{1.1}{1}\right) + \left(\ln 1 + \frac{1.0}{1}\right) + \left(\ln 1 + \frac{8.0}{1}\right)\right] \\
  &= \frac{1}{5}(1.2 + 0.9 + 1.1 + 1.0 + 8.0) = 2.440 \\[4pt]
\QLIKE_B &= \frac{1}{5}\bigl[(\ln 1.3 + 1.2/1.3) + (\ln 0.8 + 0.9/0.8) \\
  &\qquad + (\ln 1.0 + 1.1/1.0) + (\ln 1.1 + 1.0/1.1) + (\ln 3.0 + 8.0/3.0)\bigr] \\
  &= \frac{1}{5}(1.185 + 0.902 + 1.100 + 1.004 + 3.766) = 1.591
\end{align*}

QLIKE also ranks Model B first here (1.591 < 2.440), but the gap is much smaller: Model B wins by 35\% under QLIKE versus 49\% under MSE.
In samples where the crisis day is less extreme, MSE and QLIKE rankings can reverse entirely.
\end{workedexample}

\begin{keyidea}[Always Report QLIKE as Primary]
Use $\QLIKE$ as your primary loss function for volatility forecast evaluation.
Report MSE as a secondary check.
If the two metrics disagree on model rankings, the QLIKE ranking is more reliable for practical forecasting because it is less distorted by a few extreme days.
\end{keyidea}

\begin{projectconnection}[Why This Matters]
QLIKE is THE primary evaluation metric for your project.
When you report that your ML model beats HAR, the headline number is the percentage reduction in QLIKE.
The asymmetry is critical: QLIKE penalizes you more for underestimating vol than overestimating it (through the $\sigma^2_t / h_t$ ratio), which aligns with risk management priorities where underestimating vol means holding too much risk.
Target a 30--80 bps QLIKE improvement over HAR to claim a meaningful result.
\end{projectconnection}


% ══════════════════════════════════════════════════════════════
\subsection{Retransformation Bias}
\label{sec:eval-retransformation}
% ══════════════════════════════════════════════════════════════

Many volatility models forecast in log space because $\log \RV_t$ is more Gaussian, more homoskedastic, and better behaved for regression than raw $\RV_t$.
The HAR-log model, for example, regresses $\log \RV_{t+1}$ on lagged log realized variances.
But when you need a level-space forecast (e.g., for portfolio variance targeting or VaR computation), you must exponentiate the log forecast back to levels.
This innocent-looking step introduces a systematic downward bias known as \textbf{retransformation bias}.

\subsubsection{The Problem: Jensen's Inequality}

The root cause is \textbf{Jensen's inequality}: for any convex function $g$ and non-degenerate random variable $X$,
\begin{equation}
\label{eq:jensen}
\E\bigl[g(X)\bigr] > g\bigl(\E[X]\bigr).
\end{equation}
The exponential function is convex, so $\E[\exp(X)] > \exp(\E[X])$.

Suppose your log-space model produces a point forecast $\widehat{\log \RV}_{t+1}$.
The naive level-space forecast is:
\[
\widehat{\RV}^{\text{naive}}_{t+1} = \exp\!\bigl(\widehat{\log \RV}_{t+1}\bigr).
\]
But because the log-space forecast has estimation error, the true conditional expectation of $\RV_{t+1}$ is \emph{larger} than this.
Exponentiating the conditional mean of the log gives you something systematically below the conditional mean of the level.
Every single forecast is biased low.

\begin{intuition}[In Plain English]
Think of it this way.
Your log-space model says ``the average of $\log \RV$ tomorrow is 0.5.''
But the average of $\RV$ tomorrow is not $\exp(0.5) = 1.65$.
It is higher, because the distribution of $\log \RV$ has spread around 0.5, and the exponential function amplifies high values more than it shrinks low values.
The more uncertain your log-space forecast (wider error distribution), the larger the gap between $\exp(\E[\log \RV])$ and $\E[\RV]$.
\end{intuition}

\subsubsection{The Correction Formula}

If log-space forecast errors are approximately Gaussian with variance $\hat{\sigma}^2_\varepsilon$, the bias-corrected level-space forecast is:
\begin{equation}
\label{eq:retransform}
\widehat{\RV}_{t+1} = \exp\!\left(\widehat{\log \RV}_{t+1} + \frac{\hat{\sigma}^2_\varepsilon}{2}\right)
\end{equation}
where:
\begin{itemize}[nosep]
  \item $\widehat{\log \RV}_{t+1}$ is the log-space point forecast,
  \item $\hat{\sigma}^2_\varepsilon$ is the estimated variance of the log-space forecast errors $\varepsilon_t = \log \RV_t - \widehat{\log \RV}_t$,
  \item the $\hat{\sigma}^2_\varepsilon / 2$ term is the correction that offsets Jensen's inequality.
\end{itemize}

This formula comes from the moment generating function of a Gaussian: if $\varepsilon \sim \N(0, \sigma^2_\varepsilon)$, then $\E[\exp(\varepsilon)] = \exp(\sigma^2_\varepsilon / 2)$.
The corrected forecast multiplies the naive exponential by this factor.

\begin{intuition}[In Plain English]
The correction adds half the forecast error variance back before exponentiating.
It says: ``My best guess for $\log \RV$ is $\hat{y}$, but there is uncertainty of $\hat{\sigma}^2_\varepsilon$ around that guess.
Because exponentiation amplifies upside errors more than downside errors, I need to nudge my forecast upward by $\hat{\sigma}^2_\varepsilon / 2$ to get the right average in level space.''
When forecast errors are small ($\hat{\sigma}^2_\varepsilon \approx 0$), the correction is negligible.
When they are large, as in long-horizon forecasts or during volatile regimes, it can be substantial.
\end{intuition}

\subsubsection{How Large Is the Bias?}

\begin{workedexample}{Retransformation Bias Magnitude}
Suppose your HAR-log model has log-space forecast error variance $\hat{\sigma}^2_\varepsilon = 0.20$ (a realistic value for 5-day-ahead forecasts of log daily $\RV$ on equity indices).

\textbf{Correction factor:}
\[
\exp\!\left(\frac{0.20}{2}\right) = \exp(0.10) \approx 1.105
\]

This means the naive forecast underestimates the true conditional mean by about 10.5\%.
On a day where $\widehat{\log \RV}_{t+1} = -4.0$ (corresponding to annualized vol of about 14\%), the naive forecast is $\exp(-4.0) = 0.0183$, while the corrected forecast is $0.0183 \times 1.105 = 0.0202$.

For 1-day-ahead forecasts with $\hat{\sigma}^2_\varepsilon \approx 0.08$, the correction factor is $\exp(0.04) \approx 1.04$, a 4\% adjustment.
For 22-day-ahead forecasts with $\hat{\sigma}^2_\varepsilon \approx 0.35$, it reaches $\exp(0.175) \approx 1.19$, nearly a 20\% adjustment.
\end{workedexample}

\begin{warning}[Without Correction, Every Forecast Is Biased Low]
If you forecast in log space and naively exponentiate, your Mincer--Zarnowitz regression (Section~\ref{sec:eval-mz}) will show $a > 0$ (systematic under-prediction) and the bias grows with forecast uncertainty.
During high-volatility regimes, when $\hat{\sigma}^2_\varepsilon$ is largest, the bias is at its worst, precisely when accurate forecasts matter most for risk management.
\end{warning}

\begin{keyidea}[Estimating the Correction Variance]
The correction requires $\hat{\sigma}^2_\varepsilon$, the variance of log-space forecast errors.
In practice, estimate this from the training sample or from a rolling window of recent out-of-sample errors.
Using a rolling window (e.g., the trailing 60 trading days) allows the correction to adapt to regime changes: during calm markets $\hat{\sigma}^2_\varepsilon$ is small and the correction is minor; during crises it grows, appropriately increasing the level-space forecast.
\end{keyidea}

\begin{projectconnection}[Why This Matters]
Your project forecasts $\log \RV_{t+1}$ (because log realized variance is better behaved for HAR and LightGBM regressions), but $\QLIKE$ evaluation and downstream applications (volatility targeting, VaR) require level-space forecasts.
Apply the retransformation correction from Equation~\eqref{eq:retransform} whenever you convert back to levels.
Without it, you introduce a systematic negative bias that inflates $\QLIKE$ loss and makes your MZ regression show $a > 0$.
The correction is trivially cheap to compute, so there is no reason to skip it \citep{Patton2011}.
\end{projectconnection}


% ── Diagram: QLIKE vs MSE Penalty Asymmetry ──
\begin{figure}[ht]
\centering
\begin{tikzpicture}
\begin{axis}[
  width=11cm, height=7cm,
  xlabel={Forecast ratio $h_t / \sigma^2_t$},
  ylabel={Loss contribution (normalized)},
  xmin=0.2, xmax=3.5,
  ymin=0, ymax=12,
  legend pos=north west,
  legend style={font=\small},
  grid=major,
  grid style={gray!20},
  every axis label/.style={font=\small},
  tick label style={font=\footnotesize},
]

% MSE: (sigma^2 - h)^2 / sigma^4 = (1 - h/sigma^2)^2 * sigma^4 / sigma^4
% Normalized: let sigma^2 = 1, then MSE = (1 - r)^2 where r = h/sigma^2
\addplot[thick, blue, domain=0.2:3.5, samples=100]
  {(1 - x)^2};
\addlegendentry{MSE: $(1 - h_t/\sigma^2_t)^2$}

% QLIKE: ln(h) + sigma^2/h, normalized by subtracting minimum (at h=sigma^2)
% min of ln(r) + 1/r is at r=1, value = 0 + 1 = 1
% So QLIKE contribution = ln(r) + 1/r - 1 (shifted so min = 0)
\addplot[thick, red, domain=0.2:3.5, samples=100]
  {ln(x) + 1/x - 1};
\addlegendentry{QLIKE: $\ln(h_t/\sigma^2_t) + \sigma^2_t/h_t - 1$}

% Vertical line at perfect forecast
\addplot[dashed, black, thin] coordinates {(1,0) (1,12)};
\node[font=\scriptsize, anchor=south] at (axis cs:1,10.5) {Perfect forecast};

% Annotations for asymmetry
\draw[-{Stealth}, thick, red!70!black] (axis cs:0.35,7) -- (axis cs:0.35,4.5);
\node[font=\scriptsize, text width=2.2cm, align=center, red!70!black] at (axis cs:0.55,8.5) {QLIKE: harsh penalty for under-prediction};

\draw[-{Stealth}, thick, blue!70!black] (axis cs:2.8,4.5) -- (axis cs:2.8,3.2);
\node[font=\scriptsize, text width=2.5cm, align=center, blue!70!black] at (axis cs:2.8,5.8) {MSE: symmetric penalty};

\end{axis}
\end{tikzpicture}
\caption{QLIKE vs.\ MSE penalty as a function of the forecast ratio $h_t / \sigma^2_t$.
Both losses are minimized at the perfect forecast ($h_t = \sigma^2_t$, ratio $= 1$).
MSE penalizes over- and under-prediction symmetrically.
QLIKE penalizes under-prediction (ratio $< 1$) much more harshly than over-prediction (ratio $> 1$), matching the asymmetric risk preferences in volatility forecasting: underestimating vol means holding too much risk.}
\label{fig:qlike-vs-mse}
\end{figure}


% ══════════════════════════════════════════════════════════════
\section{Mincer--Zarnowitz Regressions}
\label{sec:eval-mz}
% ══════════════════════════════════════════════════════════════

$\QLIKE$ tells you which model has lower average loss, but it does not tell you \emph{why} a forecast is bad.
The Mincer--Zarnowitz regression is a simple diagnostic that decomposes forecast errors into bias and inefficiency.

\subsection{The Regression}

Regress realized variance on the forecast:
\begin{equation}
\label{eq:mz}
\sigma^2_t = a + b \cdot h_t + \varepsilon_t
\end{equation}
where:
\begin{itemize}[nosep]
  \item $\sigma^2_t$ is realized variance (left-hand side),
  \item $h_t$ is the forecast (right-hand side),
  \item $a$ is the intercept (bias term),
  \item $b$ is the slope (efficiency term),
  \item $\varepsilon_t$ is the residual.
\end{itemize}

\begin{intuition}[In Plain English]
The Mincer--Zarnowitz regression asks: ``If I plot realized variance against my forecast, do the points lie along the 45-degree line?''
A perfect forecast gives intercept $a = 0$ (no constant bias) and slope $b = 1$ (every unit increase in the forecast corresponds to exactly one unit increase in reality).
If $b < 1$, your forecast is too timid; if $b > 1$, it overreacts.
\end{intuition}

\begin{projectconnection}[Why This Matters]
After fitting your ML model, run the MZ regression before anything else.
HAR models typically show $b$ slightly below 1 (they smooth too much in high-vol regimes), and your ML extension should fix this.
If your HARQ-ML model still shows $b = 0.85$, you know the improvement needs to come from making the forecast more reactive to recent variance spikes, not from reducing average bias.
\end{projectconnection}

\begin{definition}[Unbiased and Efficient Forecast]
A forecast is \textbf{unbiased} if $a = 0$ (no systematic over- or under-prediction) and \textbf{efficient} if $b = 1$ (the forecast captures the full scale of variance movements).
Test the joint hypothesis $H_0: a = 0, \, b = 1$ with a standard F-test \citep{MincerZarnowitz1969}.
\end{definition}

\subsection{Interpreting Deviations}

\begin{itemize}[nosep]
  \item $a > 0$, $b \approx 1$: the forecast systematically under-predicts by a constant.
  \item $a \approx 0$, $b < 1$: the forecast under-reacts to variance movements (too smooth).
  \item $a \approx 0$, $b > 1$: the forecast over-reacts (too volatile).
  \item $R^2$: the fraction of realized variance variation explained by the forecast. Higher is better.
\end{itemize}

\begin{intuition}[Mincer--Zarnowitz as a Diagnostic]
Think of the MZ regression as a ``bias and calibration check.''
$\QLIKE$ tells you the overall score; MZ tells you what to fix.
If $b = 0.7$, your forecast is too smooth: it needs to react more aggressively to recent information.
If $a = 0.003$, your forecast systematically under-predicts by about 0.3 variance points.
\end{intuition}

\begin{warning}[Use HAC Standard Errors]
Volatility forecast errors are serially correlated (today's error predicts tomorrow's error) because volatility clusters (Chapter~\ref{ch:garch}).
Use Newey--West (HAC) standard errors in the MZ regression.
OLS standard errors will be too small, leading you to reject $H_0$ too often.
\end{warning}


% ══════════════════════════════════════════════════════════════
\section{The Diebold--Mariano Test}
\label{sec:eval-dm}
% ══════════════════════════════════════════════════════════════

You now have a loss function ($\QLIKE$) and a diagnostic (MZ regression).
The next question is: given two models, is the difference in loss \emph{statistically significant}, or could it be sampling noise?

\subsection{Setup}

Suppose you have two volatility forecasts, $h^A_t$ and $h^B_t$, and a loss function $L$ (use $\QLIKE$).
Define the \textbf{loss differential}:
\begin{equation}
\label{eq:dm-diff}
d_t = L(\sigma^2_t, h^A_t) - L(\sigma^2_t, h^B_t)
\end{equation}
where:
\begin{itemize}[nosep]
  \item $d_t$ is the difference in loss on day $t$,
  \item $L(\sigma^2_t, h^A_t)$ is the loss of model A on day $t$,
  \item $L(\sigma^2_t, h^B_t)$ is the loss of model B on day $t$.
\end{itemize}

If $d_t > 0$ on average, model B has lower loss (model B wins).
The question is whether $\bar{d}$ is significantly different from zero.

\begin{intuition}[In Plain English]
The loss differential $d_t$ is simply the daily scorecard: on each day, which model had a worse QLIKE score?
Some days model A wins, some days model B wins.
You are asking whether model B wins often enough, by enough, that it cannot be explained by random chance.
\end{intuition}

\subsection{The Test Statistic}

\begin{equation}
\label{eq:dm-stat}
\text{DM} = \frac{\bar{d}}{\sqrt{\widehat{\text{Var}}(\bar{d})}}
\end{equation}
where:
\begin{itemize}[nosep]
  \item $\bar{d} = \frac{1}{T}\sum_{t=1}^T d_t$ is the mean loss differential,
  \item $\widehat{\text{Var}}(\bar{d})$ is estimated using HAC (Newey--West) standard errors to account for serial correlation in $d_t$,
  \item Under $H_0: \E[d_t] = 0$, the DM statistic is asymptotically standard normal \citep{DieboldMariano1995}.
\end{itemize}

\begin{intuition}[In Plain English]
The DM statistic is a t-test on the average loss difference.
The numerator is ``how much better is model B on average?'' and the denominator is ``how uncertain are we about that average, given that consecutive days are correlated?''
If the ratio exceeds roughly 2, you have a statistically significant winner at the 5\% level.
\end{intuition}

\begin{projectconnection}[Why This Matters]
The DM test is the formal statistical test you need to claim ``my ML model significantly beats HAR.''
Without it, a reviewer can dismiss any QLIKE improvement as sampling noise.
When you report results, the DM $p$-value goes right next to the QLIKE numbers.
If $p > 0.05$, your improvement is not credible regardless of how good the point estimate looks.
\end{projectconnection}

\begin{prereq}[HAC Standard Errors]
When observations are serially correlated, the usual variance estimator $\widehat{\text{Var}}(\bar{d}) = s^2_d / T$ is biased downward.
\textbf{Heteroskedasticity and Autocorrelation Consistent (HAC)} estimators, such as Newey--West, correct for this by including autocovariances up to some lag $\ell$.
A common rule of thumb is $\ell = \lfloor T^{1/3} \rfloor$.
For $T = 1{,}000$ days, this gives $\ell = 10$.
\end{prereq}

\subsection{Worked Example}

\begin{workedexample}{Diebold--Mariano Test}
You evaluate HAR and LightGBM forecasts over $T = 500$ trading days using $\QLIKE$ loss.

\textbf{Given:}
\begin{itemize}[nosep]
  \item Mean QLIKE loss differential: $\bar{d} = 0.023$ (HAR has higher loss, so LightGBM wins on average).
  \item HAC standard error of $\bar{d}$: $\text{SE}_{\text{HAC}} = 0.011$.
\end{itemize}

\textbf{Compute:}
\[
\text{DM} = \frac{0.023}{0.011} = 2.09
\]

\textbf{Interpret:}
Under $H_0$, $\text{DM} \sim \N(0,1)$.
The two-sided $p$-value is $2 \times \Phi(-2.09) \approx 0.037$.
At the 5\% level, you reject $H_0$: the LightGBM improvement over HAR is statistically significant.

At the 1\% level, you would not reject.
Report the $p$-value and let the reader decide.
\end{workedexample}

\begin{warning}[Small-Sample Correction]
\citet{DieboldMariano1995} derived the test for large samples.
With fewer than 100 observations, use the modified DM statistic from \citet{HarveyLeybourneNewbold1997}, which uses a $t$-distribution with $T-1$ degrees of freedom and applies a finite-sample correction factor.
\end{warning}


% ══════════════════════════════════════════════════════════════
\section{The Model Confidence Set}
\label{sec:eval-mcs}
% ══════════════════════════════════════════════════════════════

The Diebold--Mariano test compares models in pairs.
With $M$ models, you would need $\binom{M}{2}$ pairwise tests, and the more tests you run, the more likely you are to find a ``significant'' difference by chance.
The Model Confidence Set solves this by comparing all models simultaneously.

\subsection{Intuition}

\begin{intuition}[The Model Confidence Set as a Tournament]
Imagine a round-robin tournament.
Instead of declaring a single winner, the MCS procedure eliminates models that are \emph{significantly worse} than the others and returns the set of survivors.
The survivors are statistically indistinguishable from each other at the chosen confidence level.
\end{intuition}

\subsection{The Procedure}

The MCS algorithm of \citet{HansenLundeNason2011} works as follows:

\begin{enumerate}[nosep]
  \item Start with the full set of $M$ models: $\mathcal{M}_0 = \{1, 2, \ldots, M\}$.
  \item Test the null hypothesis $H_0$: all models in the current set have equal expected loss.
  \item If $H_0$ is rejected at significance level $\alpha$, identify and remove the worst model (the one with the highest average loss).
  \item Repeat steps 2--3 on the reduced set until $H_0$ is not rejected.
  \item The surviving set $\widehat{\mathcal{M}}^*_\alpha$ is the \textbf{Model Confidence Set} at level $\alpha$.
\end{enumerate}

\begin{definition}[Model Confidence Set]
The Model Confidence Set $\widehat{\mathcal{M}}^*_\alpha$ at significance level $\alpha$ contains all models whose forecasting performance is not significantly worse than the best model.
Formally, it satisfies:
\[
\Pr\left(\mathcal{M}^* \subseteq \widehat{\mathcal{M}}^*_\alpha\right) \geq 1 - \alpha
\]
where $\mathcal{M}^*$ is the (unknown) set of truly best models.
\end{definition}

\begin{keyresult}[Hansen, Lunde, and Nason (2011): The Gold Standard for Multi-Model Comparison]
\citet{HansenLundeNason2011} show that the MCS controls the familywise error rate: the probability of incorrectly excluding any truly best model is at most $\alpha$.
The MCS produces a \emph{set}, not a ranking.
Reporting ``these 4 models are in the 90\% MCS'' is more honest than reporting ``model X has the lowest QLIKE'' when differences are small.
\end{keyresult}

% ── Diagram: Model Confidence Set ──
\begin{figure}[ht]
\centering
\begin{tikzpicture}[
  every node/.style={font=\small},
]
% Outer rectangle: all candidate models
\draw[thick, rounded corners=8pt] (-4.5,-2.5) rectangle (4.5,2.5);
\node[font=\small\bfseries, anchor=north west] at (-4.3,2.3) {All candidate models $\mathcal{M}_0$};

% MCS ellipse (90%)
\draw[thick, fill=intgreen!12, rounded corners=5pt] (-3.2,-1.6) rectangle (2.8,1.6);
\node[font=\footnotesize\bfseries, text=intgreen!70!black, anchor=north west] at (-3.0,1.4) {MCS$_{90\%}$};

% Models inside MCS
\node[fill=blue!15, rounded corners=3pt, minimum width=1.8cm, minimum height=0.55cm] at (-1.5,0.7) {LightGBM};
\node[fill=blue!15, rounded corners=3pt, minimum width=1.8cm, minimum height=0.55cm] at (1.2,0.7) {HAR};
\node[fill=blue!15, rounded corners=3pt, minimum width=1.8cm, minimum height=0.55cm] at (-1.5,-0.1) {GARCH};
\node[fill=blue!15, rounded corners=3pt, minimum width=1.8cm, minimum height=0.55cm] at (1.2,-0.1) {Rand.\ Forest};

% Models outside MCS (eliminated)
\node[fill=warnred!15, rounded corners=3pt, minimum width=1.8cm, minimum height=0.55cm, text=warnred!70!black] at (3.5,-1.8) {LSTM};
\node[fill=warnred!15, rounded corners=3pt, minimum width=1.8cm, minimum height=0.55cm, text=warnred!70!black] at (3.5,2.0) {Hist.\ avg.};

% Annotation
\node[font=\scriptsize, text width=2.8cm, align=center, anchor=north] at (-0.2,-1.0) {Statistically indistinguishable\\at 90\% confidence};

% Elimination arrows
\draw[-{Stealth}, warnred, thick, dashed] (2.8,1.8) -- (3.2,1.9);
\draw[-{Stealth}, warnred, thick, dashed] (2.8,-1.5) -- (3.2,-1.7);
\node[font=\scriptsize, text=warnred!70!black, anchor=west] at (3.3,-0.5) {Eliminated};

\end{tikzpicture}
\caption{The Model Confidence Set retains all models whose performance is statistically indistinguishable from the best (green region) and eliminates models that are significantly worse (red, outside).
The surviving models cannot be ranked among themselves with statistical confidence.}
\label{fig:mcs-venn}
\end{figure}

\subsection{Practical Use}

Report which models survive at both $\alpha = 0.10$ and $\alpha = 0.05$:

\begin{center}
\begin{tabular}{lccc}
\toprule
Model & QLIKE & MCS $p$-value & In MCS$_{90\%}$? \\
\midrule
LightGBM + HAR features & 1.423 & 1.000 & Yes \\
HAR (daily, weekly, monthly) & 1.441 & 0.482 & Yes \\
GARCH(1,1)                   & 1.467 & 0.312 & Yes \\
Random Forest                 & 1.439 & 0.551 & Yes \\
LSTM                          & 1.502 & 0.043 & No  \\
Historical average            & 1.589 & 0.001 & No  \\
\bottomrule
\end{tabular}
\end{center}

The MCS $p$-value for each model is the smallest $\alpha$ at which that model would be excluded.
A model with MCS $p$-value $> 0.10$ survives in the 90\% MCS.
In this example, four models are statistically indistinguishable at the 90\% level; the LSTM and historical average are significantly worse.

\begin{keyidea}[MCS as a Humility Device]
The MCS often reveals an uncomfortable truth: many models that look different in QLIKE are statistically indistinguishable.
A 2\% QLIKE improvement is rarely significant with 3--5 years of daily data.
If your fancy model is in the same MCS as HAR, be honest about it.
\end{keyidea}

\begin{projectconnection}[Why This Matters]
Your model needs to be IN the Model Confidence Set, and ideally, simpler baselines like raw HAR should be excluded.
If your LightGBM model and plain HAR both survive in the 90\% MCS, you cannot honestly claim superiority.
The MCS is also your defense: if a reviewer asks ``why not use an LSTM?'', you can show it was eliminated from the MCS.
Use the \texttt{MCS} package in R or the \texttt{arch} library in Python to compute MCS $p$-values.
\end{projectconnection}


% ══════════════════════════════════════════════════════════════
\section{Purged K-Fold Cross-Validation with Embargo}
\label{sec:eval-purgedcv}
% ══════════════════════════════════════════════════════════════

The tests above (DM, MCS) evaluate forecasts on a held-out sample.
But how do you \emph{select} the model and tune hyperparameters in the first place?
Standard K-fold cross-validation fails catastrophically on time series data.
This section explains why and introduces the fix.

\subsection{Why Standard K-Fold Fails}

\begin{prereq}[K-Fold Cross-Validation]
In standard K-fold CV, you split data into $K$ equally sized folds, train on $K-1$ folds, test on the remaining fold, and rotate.
This works when observations are independent (e.g., images, text documents).
It does \emph{not} work when observations are serially dependent.
\end{prereq}

Consider 5-fold CV on 1,250 trading days (5 years).
Fold 1 = days 1--250, fold 2 = days 251--500, and so on.
When testing on fold 2 (days 251--500), you train on folds 1, 3, 4, 5 (days 1--250 \emph{and} 501--1250).

The problem: volatility on day 501 is highly correlated with volatility on day 500 (the last day of the test set).
Training on day 501 while testing on day 500 is using the future to predict the past.
Worse, if your labels use multi-day returns (e.g., 5-day forward realized variance), then the label for day 498 overlaps with the label for day 502; the training and test sets share information through the label construction.

% ── Diagram: Purged K-Fold Timeline ──
\begin{figure}[ht]
\centering
\begin{tikzpicture}[
  x=0.018cm,
  fold/.style={minimum height=0.7cm, anchor=west, font=\small\bfseries},
  purge/.style={minimum height=0.7cm, anchor=west, font=\small, pattern=north east lines, pattern color=warnred!60},
  embargo/.style={minimum height=0.7cm, anchor=west, font=\small, pattern=dots, pattern color=prereqpurple!80}
]

% Timeline axis
\draw[thick, -{Stealth}] (0,0) -- (680,0);
\foreach \x/\lab in {0/0, 125/250, 250/500, 375/750, 500/1000, 625/1250} {
  \draw (\x, -0.15) -- (\x, 0.15);
  \node[below, font=\footnotesize] at (\x, -0.2) {\lab};
}
\node[below, font=\small] at (312, -0.8) {Trading days};

% Row 1: Standard fold assignment
\node[font=\small\bfseries, anchor=east] at (-10, 1.5) {Folds:};
\fill[blue!15] (0, 1.1) rectangle (125, 1.9);
\node[font=\small] at (62.5, 1.5) {Fold 1};
\fill[warnred!25] (125, 1.1) rectangle (250, 1.9);
\node[font=\small\bfseries, text=warnred] at (187.5, 1.5) {Test};
\fill[blue!15] (250, 1.1) rectangle (375, 1.9);
\node[font=\small] at (312.5, 1.5) {Fold 3};
\fill[blue!15] (375, 1.1) rectangle (500, 1.9);
\node[font=\small] at (437.5, 1.5) {Fold 4};
\fill[blue!15] (500, 1.1) rectangle (625, 1.9);
\node[font=\small] at (562.5, 1.5) {Fold 5};

% Row 2: With purging and embargo
\node[font=\small\bfseries, anchor=east] at (-10, 3.2) {Purged:};
\fill[blue!15] (0, 2.8) rectangle (100, 3.6);
\node[font=\small] at (50, 3.2) {Train};
\fill[warnred!15, draw=warnred, thick, densely dashed] (100, 2.8) rectangle (125, 3.6);
\node[font=\tiny, text=warnred, rotate=90] at (112.5, 3.2) {purge};
\fill[warnred!25] (125, 2.8) rectangle (250, 3.6);
\node[font=\small\bfseries, text=warnred] at (187.5, 3.2) {Test};
\fill[prereqpurple!15, draw=prereqpurple, thick, densely dashed] (250, 2.8) rectangle (275, 3.6);
\node[font=\tiny, text=prereqpurple, rotate=90] at (262.5, 3.2) {embargo};
\fill[blue!15] (275, 2.8) rectangle (375, 3.6);
\node[font=\small] at (325, 3.2) {Train};
\fill[blue!15] (375, 2.8) rectangle (500, 3.6);
\node[font=\small] at (437.5, 3.2) {Train};
\fill[blue!15] (500, 2.8) rectangle (625, 3.6);
\node[font=\small] at (562.5, 3.2) {Train};

% Annotations
\draw[warnred, thick, <->] (100, 4.0) -- (125, 4.0);
\node[above, font=\scriptsize, text=warnred] at (112.5, 4.0) {25 days};
\draw[prereqpurple, thick, <->] (250, 4.0) -- (275, 4.0);
\node[above, font=\scriptsize, text=prereqpurple] at (262.5, 4.0) {25 days};

\end{tikzpicture}
\caption{Purged K-fold CV with embargo ($K=5$, $T=1{,}250$, embargo $= 2\%$).
\textbf{Top row:} standard fold assignment with fold 2 as the test set.
\textbf{Bottom row:} after purging and embargo.
The 25 days before the test set (\textcolor{warnred}{purge zone}) are removed from training to prevent label overlap.
The 25 days after the test set (\textcolor{prereqpurple}{embargo zone}) are removed to prevent information leakage from serial correlation.
Training uses only the remaining blue regions.}
\label{fig:purged-cv}
\end{figure}

\subsection{The Fix: Purging and Embargo}

\citet{LopezdePrado2018} introduces two modifications to standard K-fold CV:

\begin{definition}[Purging]
\textbf{Purging} removes from the training set any observations whose label windows overlap with the test period.
If labels are constructed from $\tau$-day forward returns, remove training observations within $\tau$ days before the start of the test fold.
\end{definition}

\begin{definition}[Embargo]
\textbf{Embargo} removes an additional buffer of training observations \emph{after} the end of the test fold.
This guards against serial correlation in features: day $t+1$ features are correlated with day $t$ features, so training on day $t+1$ while testing on day $t$ leaks information.
A typical embargo is 1--2\% of total sample size.
\end{definition}

\subsection{Worked Example}

\begin{workedexample}{Purged 5-Fold CV}
\textbf{Setup:} $T = 1{,}250$ daily observations (5 years), $K = 5$ folds, labels use 5-day forward realized variance, embargo $= 2\%$ of $T = 25$ days.

\textbf{Fold 2 as test set} (days 251--500):
\begin{enumerate}[nosep]
  \item \textbf{Standard training set:} days 1--250 and 501--1250 (1,000 days).
  \item \textbf{After purging:} remove days 246--250 from training (their 5-day labels overlap with the test period). Training loses 5 days.
  \item \textbf{After embargo:} remove days 501--525 from training (serial correlation buffer). Training loses 25 more days.
  \item \textbf{Final training set:} days 1--245 and 526--1250 (970 days).
\end{enumerate}

You lose 30 training observations per fold (3\% of the total). This is a small price for preventing look-ahead bias.
\end{workedexample}

\begin{projectconnection}[Why This Matters]
Your vol forecasting labels use multi-day forward realized variance (typically 1-day or 5-day), which means label windows overlap across consecutive observations.
Standard K-fold would train on day 502 (whose label includes returns from days 502--506) while testing on day 500 (whose label includes days 500--504).
The overlapping days 502--504 leak test information into training.
Purging removes this overlap; embargo handles the residual serial correlation in features like lagged RV.
Use \texttt{sklearn.model\_selection.TimeSeriesSplit} as a starting point, then add purging and embargo manually or use the \texttt{purged\_cv} implementation from \citet{LopezdePrado2018}.
\end{projectconnection}

\begin{warning}[Random K-Fold on Time Series Is Catastrophic]
Random K-fold on time series data (shuffling observations before splitting) is the single most common evaluation error in ML-for-finance papers.
A model trained on January and March, tested on February, has literally seen the future.
Reported accuracy will be dramatically inflated; out-of-sample performance will collapse.
Always use purged CV, expanding-window, or walk-forward evaluation for time series.
\end{warning}


% ══════════════════════════════════════════════════════════════
\section{The Deflated Sharpe Ratio}
\label{sec:eval-dsr}
% ══════════════════════════════════════════════════════════════

Everything above evaluates \emph{forecasts} of volatility.
But volatility forecasts are often embedded in trading strategies (volatility targeting, variance risk premium trading; see Chapters~\ref{ch:vrp} and~\ref{ch:applications}).
The standard performance metric for strategies is the Sharpe ratio.
This section explains why raw Sharpe ratios are misleading when you have tried multiple strategies, and how to correct them.

\subsection{The Multiple Testing Problem}

\begin{intuition}[Sharpe Ratios and Coin Flips]
Suppose you flip 30 fair coins 250 times each and report only the coin with the most heads.
That coin will show a ``success rate'' well above 50\%, but it has no skill; you simply selected the luckiest coin.
The same logic applies to Sharpe ratios: if you try 30 feature sets and report the best one, the expected maximum Sharpe ratio under the null (no skill) is not zero.
\end{intuition}

\citet{Bailey2014DSR} derive the expected maximum Sharpe ratio under the null when $N$ independent strategies are tested:
\begin{equation}
\label{eq:expected-max-sr}
\E\bigl[\max_{i=1,\ldots,N} \SR_i\bigr] \approx \sqrt{2 \ln N}
\end{equation}
where:
\begin{itemize}[nosep]
  \item $N$ is the number of independent strategies (or feature sets, or hyperparameter combinations) tested,
  \item $\SR_i$ is the Sharpe ratio of strategy $i$ under the null (all strategies have true $\SR = 0$),
  \item The approximation comes from extreme value theory for Gaussian maxima.
\end{itemize}

For $N = 30$, this gives $\E[\max \SR] \approx \sqrt{2 \ln 30} \approx 2.61$.
A reported Sharpe of 1.5 after 30 trials is \emph{below} what you would expect from pure luck.

\begin{intuition}[In Plain English]
This formula says: the more strategies you try, the higher the Sharpe ratio you should expect from the luckiest one, even if none of them have any real skill.
It grows slowly (as $\sqrt{\ln N}$), but 30 trials already pushes the luck threshold to a Sharpe of 2.6.
Your observed Sharpe must exceed this threshold to be credible.
\end{intuition}

\begin{projectconnection}[Why This Matters]
If you test 20 hyperparameter configurations for your vol-timing strategy, the expected maximum Sharpe under pure luck is $\sqrt{2 \ln 20} \approx 2.45$.
Any backtest Sharpe below this number is entirely consistent with having no skill.
This is why you must log every experiment from the start: $N$ only grows, and forgetting trials inflates your apparent performance.
\end{projectconnection}

\subsection{The DSR Formula}

The Deflated Sharpe Ratio adjusts the observed Sharpe ratio for the number of trials:
\begin{equation}
\label{eq:dsr}
\DSR = \Phi\!\left(\frac{(\widehat{\SR} - \SR_0)\sqrt{T-1}}{\sqrt{1 - \hat{\gamma}_3 \widehat{\SR} + \frac{\hat{\gamma}_4 - 1}{4}\widehat{\SR}^2}}\right)
\end{equation}
where:
\begin{itemize}[nosep]
  \item $\DSR \in [0, 1]$ is the probability that the observed Sharpe exceeds the multiple-testing threshold (higher is better),
  \item $\widehat{\SR}$ is the observed (annualized) Sharpe ratio of the best strategy,
  \item $\SR_0 = \sqrt{2 \ln N}$ is the expected maximum Sharpe under the null, with $N =$ number of trials,
  \item $T$ is the number of return observations,
  \item $\hat{\gamma}_3$ is the sample skewness of returns,
  \item $\hat{\gamma}_4$ is the sample kurtosis of returns,
  \item $\Phi(\cdot)$ is the standard normal CDF.
\end{itemize}

\begin{intuition}[In Plain English]
The DSR converts your observed Sharpe ratio into a probability: ``What is the chance that this Sharpe is real, given how many strategies I tried?''
It subtracts the luck threshold ($\SR_0$) from your observed Sharpe, scales by sample size, and adjusts for skewness and fat tails in your returns.
A DSR near 1 means your Sharpe is almost certainly genuine; a DSR near 0 means it is probably luck.
\end{intuition}

\begin{projectconnection}[Why This Matters]
If your vol-forecasting project includes a variance risk premium trading strategy (Chapter~\ref{ch:vrp}), you will need to report DSR alongside the raw Sharpe.
With the typical 10--30 experiments you will run during hyperparameter tuning, even a Sharpe of 1.5 can be entirely consistent with luck.
DSR $> 0.95$ is the bar for a credible backtest result.
\end{projectconnection}

\begin{keyresult}[Bailey and Lopez de Prado (2014): The Deflated Sharpe Ratio]
\citet{Bailey2014DSR} show that ignoring the number of trials leads to systematic over-reporting of Sharpe ratios in backtested strategies.
The DSR corrects for this by benchmarking the observed Sharpe against the expected maximum under the null.
A DSR above 0.95 provides evidence that the strategy's Sharpe ratio is unlikely to have arisen from multiple testing alone.
\end{keyresult}

\subsection{Worked Example}

\begin{workedexample}{Deflated Sharpe Ratio}
You test $N = 20$ volatility-timing strategies over $T = 1{,}260$ trading days (5 years of daily returns).
The best strategy achieves $\widehat{\SR} = 1.8$ (annualized).
Returns have skewness $\hat{\gamma}_3 = -0.3$ and kurtosis $\hat{\gamma}_4 = 4.2$.

\textbf{Step 1: Compute the null benchmark.}
\[
\SR_0 = \sqrt{2 \ln 20} \approx \sqrt{5.99} \approx 2.45
\]

\textbf{Step 2: Compute the DSR numerator.}
\[
(\widehat{\SR} - \SR_0)\sqrt{T - 1} = (1.8 - 2.45)\sqrt{1259} = (-0.65)(35.48) = -23.06
\]

\textbf{Step 3: Compute the DSR denominator.}
\[
\sqrt{1 - (-0.3)(1.8) + \frac{4.2 - 1}{4}(1.8)^2} = \sqrt{1 + 0.54 + 2.592} = \sqrt{4.132} \approx 2.033
\]

\textbf{Step 4: Compute DSR.}
\[
\DSR = \Phi\!\left(\frac{-23.06}{2.033}\right) = \Phi(-11.35) \approx 0.000
\]

The DSR is essentially zero.
Despite an impressive-looking Sharpe of 1.8, the strategy does not survive correction for 20 trials.
The expected maximum Sharpe under pure luck ($\SR_0 = 2.45$) exceeds the observed Sharpe.
You cannot reject the null that this is the luckiest of 20 random strategies.
\end{workedexample}

\begin{keyidea}[Every Experiment Counts as a Trial]
The DSR requires you to know $N$, the total number of strategies tested.
This includes every feature set, every hyperparameter grid point, every ``quick look'' that influenced your final choice.
If you do not log experiments, you cannot compute an honest DSR.
This is why the experiment tracker (logging every trial) is infrastructure you build \emph{before} you start modeling.
\end{keyidea}

\begin{warning}[The Haircut Sharpe Ratio]
\citet{HarveyLiu2015} propose a complementary correction.
Rather than computing a probability, they ``haircut'' the Sharpe ratio by the amount attributable to multiple testing.
The haircut depends on the number of trials and the correlation among strategies.
Both DSR and Haircut Sharpe tell the same story: raw Sharpe ratios from backtests are systematically inflated.
Report both if possible.
\end{warning}


% ══════════════════════════════════════════════════════════════
\section{What Doesn't Work}
\label{sec:eval-pitfalls}
% ══════════════════════════════════════════════════════════════

This chapter has given you the right tools.
This section catalogs the wrong ones, so you can recognize them in other people's work and avoid them in your own.

\begin{warning}[Evaluation Pitfalls]
\begin{enumerate}[nosep]
  \item \textbf{Random K-fold on time series.}
    Shuffling observations before splitting destroys temporal structure.
    Reported accuracy is inflated; real performance collapses.
    Always use purged CV or walk-forward (Section~\ref{sec:eval-purgedcv}).

  \item \textbf{Naive out-of-sample $R^2$ without statistical tests.}
    ``Our model achieves OOS $R^2 = 3.2\%$ versus the benchmark's 2.8\%.''
    Without a DM test (Section~\ref{sec:eval-dm}) or MCS (Section~\ref{sec:eval-mcs}), you do not know whether 0.4 percentage points is signal or noise.

  \item \textbf{Training on one regime, testing on another.}
    Training on 2015--2019 (low volatility) and testing on 2020 (COVID) is not a fair evaluation; it is a regime-change stress test.
    Useful, but do not confuse it with a general OOS evaluation.

  \item \textbf{Look-ahead in feature construction.}
    Using day-$t$ VIX close to predict day-$t$ realized variance is look-ahead bias: VIX is not known until 4:15 PM, while RV accumulates throughout the day.
    All features must be known \emph{before} the forecast is made.

  \item \textbf{Reporting tiny improvements without economic significance.}
    Beating HAR by 0.5\% in QLIKE is unlikely to translate to meaningful PnL after transaction costs.
    Always pair statistical significance (DM test) with economic significance (cost-aware backtest; Chapter~\ref{ch:applications}).

  \item \textbf{Ignoring forecast variance.}
    A model that is right on average but has high forecast variance is dangerous for volatility targeting.
    Two models with identical QLIKE can differ dramatically in how ``jumpy'' their forecasts are.
    Report forecast autocorrelation and turnover alongside loss metrics.
\end{enumerate}
\end{warning}


% ══════════════════════════════════════════════════════════════
\subsection{Lookahead Bias: A Taxonomy of Four Sources}
\label{sec:eval-lookahead-taxonomy}
% ══════════════════════════════════════════════════════════════

Item 4 in the list above (look-ahead in feature construction) deserves a full subsection because lookahead bias is the single most destructive error in financial ML.
A contaminated model shows excellent in-sample performance that vanishes out of sample, wasting weeks of development time before the bug is identified.

Lookahead bias occurs whenever a feature used at prediction time contains information that would not have been available at the moment the forecast was made.
In volatility forecasting, there are four distinct sources, each with its own failure mode.
Understanding them concretely, with specific examples of how each one can leak future information into your features, is essential for building a trustworthy pipeline.


% ── Source 1: Realized Measures ──
\subsubsection{Source 1: Realized Measures}

\textbf{Realized variance} (Chapter~\ref{ch:rv}) is computed from intraday returns over a trading day.
The danger is that the boundary between ``today'' and ``tomorrow'' is not always clean.

\begin{warning}[Realized Measure Leakage]
Suppose you forecast $\RV_{t+1}$ (tomorrow's realized variance) using features that include $\RV_t$ (today's realized variance).
If you compute $\RV_t$ using returns from 9:30 AM to 4:00 PM, but the last intraday return spans 3:55--4:00 PM, that return reflects information that overlaps with the overnight period leading into day $t+1$.
In tick-level data, the problem is worse: the last trade might occur at 4:00:02 PM, technically after the close.
Even a few seconds of overlap contaminates the feature with forward-looking information.
\end{warning}

\begin{keyidea}[Prevention: Strict Temporal Cutoff]
All features for forecasting $\RV_{t+1}$ must use data from day $t$ or earlier.
Define $\RV_t$ using a fixed, consistent intraday window (e.g., 9:30 AM to 3:59 PM) and apply this cutoff uniformly across all realized measures: $\RV$, realized quarticity ($\text{RQ}$), bipower variation ($\text{BV}$), and signed components ($\RV^+$, $\RV^-$).
Timestamp every data point and assert programmatically that no feature for $\RV_{t+1}$ uses data with timestamp $> t$ close.
When in doubt, lag by one full day.
\end{keyidea}


% ── Source 2: Microstructure ──
\subsubsection{Source 2: Microstructure Features}

\textbf{Microstructure features} (Chapter~\ref{ch:noise}) are derived from limit order book (LOB) data, trade-and-quote streams, and intraday volume profiles.
They are particularly vulnerable to lookahead because they are computed over intraday windows whose boundaries must be carefully aligned with the forecast target.

\begin{warning}[LOB Feature Leakage]
Suppose you compute the \textbf{Volume-Synchronized Probability of Informed Trading (VPIN)} over the full trading day from 9:30 AM to 4:00 PM and use it as a feature for forecasting $\RV_{t+1}$.
The last hour of VPIN reflects order flow patterns driven by information that will affect the overnight return and the opening of day $t+1$.
For example, a large informed seller at 3:45 PM depresses the close and widens the spread; this information is not ``known'' to a forecaster who must act at 3:00 PM.
Full-day microstructure features effectively let you ``see'' information that accumulates between your forecast time and the close.
\end{warning}

\begin{keyidea}[Prevention: Truncate LOB Features Before Close]
Truncate all intraday microstructure features at a fixed cutoff before the market close.
A common choice is 3:00 PM (one hour before close) or even 2:00 PM for conservative pipelines.
Apply this cutoff uniformly to VPIN, Kyle's lambda, Amihud illiquidity, bid-ask spread averages, and depth imbalance measures.
Document the cutoff time as a pipeline parameter, not a magic number buried in preprocessing code.
\end{keyidea}


% ── Source 3: Options Surface ──
\subsubsection{Source 3: Options Surface}

\textbf{Implied volatility} from the options surface (Chapters~\ref{ch:garch} and~\ref{ch:applications}) is a forward-looking feature by design: it encodes the market's expectation of future volatility.
This makes it powerful but also dangerous, because the surface updates throughout the day in response to new information.

\begin{warning}[Implied Volatility Leakage]
Suppose you use the 3:30 PM VIX level as a feature for forecasting next-day $\RV$.
If a company announces earnings at 4:05 PM, the options market will begin pricing in the expected earnings move before the announcement: implied volatility rises during the last hour of trading.
The 3:30 PM VIX already reflects this anticipation.
A forecaster using this feature has access to information about the expected earnings-day volatility spike that is not ``available'' in the sense of a genuine real-time forecast made at, say, the previous close.
More subtly, the SPX implied volatility surface at 3:30 PM reflects the full day's information flow, including any macro data releases, Fed communications, or geopolitical events that occurred during the day.
\end{warning}

\begin{keyidea}[Prevention: Use Previous-Day Close or Morning Snapshot]
Use the end-of-day implied volatility surface from day $t-1$ (the previous close) or a fixed morning snapshot (e.g., 10:00 AM) as features for forecasting $\RV_{t+1}$.
Never use same-day afternoon implied volatility for next-day forecasts.
For VIX and VIX term structure features, the same rule applies: use the previous close, not the intraday value.
If you need intraday IV features for same-day forecasting (e.g., predicting afternoon $\RV$ from morning data), use a morning-only window and document it explicitly.
\end{keyidea}


% ── Source 4: Cross-Asset ──
\subsubsection{Source 4: Cross-Asset Features}

\textbf{Cross-asset features} (Chapter~\ref{ch:applications}) use data from other markets (e.g., European equities, commodities, currencies, Treasuries) to predict volatility in the target asset (e.g., SPX).
The complication is that these markets operate on different schedules, creating timestamp misalignment that can hide lookahead bias.

\begin{warning}[Cross-Asset Timing Leakage]
Suppose you use the EURO STOXX 50 realized variance as a feature for predicting SPX $\RV_{t+1}$.
European markets close at 4:30 PM CET (10:30 AM ET).
If you label the European close as ``day $t$'' data, it is available before the US close on day $t$ at 4:00 PM ET, which is fine.
But if the European data is labeled ``day $t$'' and the US target is also ``day $t$,'' you may inadvertently use European close data that overlaps with the US target window.
Worse, some cross-asset data sources (e.g., Asian markets) close well before the US open; using ``day $t$'' Asian close data for a US ``day $t{+}1$'' forecast is correct, but using it for a US ``day $t$'' forecast means the Asian data is stale and the time alignment is ambiguous.
\end{warning}

\begin{keyidea}[Prevention: Align to a Single Information Cutoff]
Define a single, global \textbf{information cutoff time} for each forecast date (e.g., previous-day US close at 4:00 PM ET).
All cross-asset features must use data from before this cutoff.
For European data, this means using the European close from the same calendar day (since it precedes the US close).
For Asian data, use the Asian close from the same calendar day (which precedes the US open).
Build an explicit timezone-aware timestamp column for every data source and assert that all feature timestamps precede the cutoff.
\end{keyidea}


% ── Summary Table ──
\subsubsection{Summary}

Table~\ref{tab:lookahead-taxonomy} collects all four sources with their specific pitfalls and prevention rules.

\begin{table}[ht]
\centering
\caption{Lookahead bias taxonomy for volatility forecasting pipelines.
Each source represents a distinct mechanism by which future information can leak into features.}
\label{tab:lookahead-taxonomy}
\small
\begin{tabular}{@{}p{2.2cm}p{5.0cm}p{5.5cm}@{}}
\toprule
\textbf{Source} & \textbf{Pitfall} & \textbf{Prevention Rule} \\
\midrule
Realized measures
  & Target-day intraday returns leak into features via boundary overlap
  & Features use data only up to day $t$ close; enforce with programmatic timestamp assertions \\
\addlinespace
Microstructure
  & Full-day LOB features (VPIN, spreads) include close-period information
  & Truncate all LOB features at a fixed cutoff (e.g., 3:00 PM) before close \\
\addlinespace
Options surface
  & Intraday IV changes reflect target-day information (e.g., earnings anticipation)
  & Use previous-day close IV or a fixed morning snapshot only \\
\addlinespace
Cross-asset
  & Mixed frequencies and timezone misalignment hide temporal overlap
  & Align all cross-asset features to a single information cutoff (previous-day US close) \\
\bottomrule
\end{tabular}
\end{table}

\begin{projectconnection}[Why This Matters]
Your pipeline ingests data from at least four different source types (tick data, LOB depth, options surface, cross-asset indices), each with its own timestamp conventions.
Build the lookahead prevention into your data pipeline as hard constraints, not as documentation that you hope developers will follow.
Specifically:
(1) add a \texttt{max\_timestamp} column to every feature table and assert it precedes the forecast cutoff;
(2) run a nightly integration test that checks no feature for $\RV_{t+1}$ uses data with timestamp $> t$ close;
(3) when adding new features, require the contributor to specify the information cutoff in the feature registry.
A single lookahead bug can invalidate months of work and is almost impossible to detect from $\QLIKE$ numbers alone: the contaminated model simply looks ``better than it should.''
\end{projectconnection}


% ══════════════════════════════════════════════════════════════
\section{Putting It All Together: An Evaluation Workflow}
\label{sec:eval-workflow}
% ══════════════════════════════════════════════════════════════

You now have all the individual tools. This section assembles them into a practical workflow you should follow for every volatility forecasting experiment.

\begin{figure}[ht]
\centering
\begin{tikzpicture}[
  node distance=1.1cm,
  every node/.style={font=\small},
  wfstep/.style={flowblock, minimum width=5cm, minimum height=0.9cm},
  decision/.style={decisionblock, minimum width=3cm, inner sep=3pt},
]
% Steps
\node[wfstep, fill=prereqpurple!10] (reserve) {1. Reserve holdout (3--6 months)};
\node[wfstep, fill=blue!8, below=of reserve] (log) {2. Initialize experiment log ($N = 0$)};
\node[wfstep, fill=blue!8, below=of log] (tune) {3. Tune with purged K-fold CV};
\node[wfstep, fill=blue!8, below=of tune] (qlike) {4. Evaluate: QLIKE (primary), MSE (secondary)};
\node[wfstep, fill=blue!8, below=of qlike] (mz) {5. MZ regression: check bias ($a=0$, $b=1$)};
\node[wfstep, fill=blue!8, below=of mz] (dm) {6. DM test: pairwise significance};
\node[wfstep, fill=keyorange!10, below=of dm] (mcs) {7. MCS: which models survive?};
\node[decision, below=1.2cm of mcs] (trade) {Strategy?};
\node[wfstep, fill=memgold!15, below left=1.2cm and -0.3cm of trade] (dsr) {8. DSR on Sharpe ratio};
\node[wfstep, fill=intgreen!10, below right=1.2cm and -0.3cm of trade] (report) {9. Report with all metrics};

% Arrows
\draw[arrow] (reserve) -- (log);
\draw[arrow] (log) -- (tune);
\draw[arrow] (tune) -- (qlike);
\draw[arrow] (qlike) -- (mz);
\draw[arrow] (mz) -- (dm);
\draw[arrow] (dm) -- (mcs);
\draw[arrow] (mcs) -- (trade);
\draw[arrow] (trade) -| node[above left, font=\scriptsize] {Yes} (dsr);
\draw[arrow] (trade) -| node[above right, font=\scriptsize] {No} (report);
\draw[arrow] (dsr) |- (report);

% Side annotation
\node[right=1.5cm of tune, font=\scriptsize, text width=3.5cm, align=left, text=prereqpurple] (note) {Every experiment increments $N$. Log feature set, hyperparameters, and QLIKE.};
\draw[-{Stealth}, prereqpurple, dashed] (note.west) -- (tune.east);

\end{tikzpicture}
\caption{Evaluation workflow for volatility forecasting.
Reserve the holdout first; log every experiment; use purged CV for tuning; evaluate with QLIKE and MZ; compare with DM and MCS; deflate the Sharpe if the forecast feeds a strategy.}
\label{fig:eval-workflow}
\end{figure}

\begin{keyidea}[The Workflow Is the Standard]
Following this workflow does not guarantee you will find a good forecast.
It guarantees that if you \emph{do} find one, the evidence will survive scrutiny.
Skip any step and a careful reader can dismiss your results.
\end{keyidea}


% ══════════════════════════════════════════════════════════════
\section{Summary}
\label{sec:eval-summary}
% ══════════════════════════════════════════════════════════════

\begin{itemize}[nosep]
  \item \textbf{MSE is proxy-robust} but over-penalizes extreme variance days, making it a poor primary metric for volatility (Section~\ref{sec:eval-mse}).

  \item \textbf{QLIKE is the preferred primary loss} for volatility forecast evaluation. It is proxy-robust \emph{and} less sensitive to outliers than MSE \citep{Patton2011} (Section~\ref{sec:eval-qlike}).

  \item \textbf{QLIKE and MSE are the only two robust losses.} Other common losses (MAE, HMSE) can reverse model rankings when the volatility proxy is noisy (Section~\ref{sec:eval-qlike}).

  \item \textbf{Retransformation bias} arises when exponentiating log-space forecasts back to levels. Apply the correction $\exp(\hat{y} + \hat{\sigma}^2_\varepsilon/2)$ to avoid systematic under-prediction that grows with forecast uncertainty \citep{Patton2011} (Section~\ref{sec:eval-retransformation}).

  \item \textbf{Mincer--Zarnowitz regressions} diagnose bias ($a \neq 0$) and inefficiency ($b \neq 1$) in forecasts. Use HAC standard errors (Section~\ref{sec:eval-mz}).

  \item \textbf{The Diebold--Mariano test} determines whether the difference in loss between two models is statistically significant, accounting for serial correlation via HAC standard errors \citep{DieboldMariano1995} (Section~\ref{sec:eval-dm}).

  \item \textbf{The Model Confidence Set} compares all models simultaneously, returning the set of models that are statistically indistinguishable from the best. It controls familywise error \citep{HansenLundeNason2011} (Section~\ref{sec:eval-mcs}).

  \item \textbf{Purged K-fold CV with embargo} prevents information leakage in time series cross-validation by removing overlapping labels (purge) and serial-correlation buffers (embargo) \citep{LopezdePrado2018} (Section~\ref{sec:eval-purgedcv}).

  \item \textbf{Random K-fold on time series is catastrophic.} It inflates reported accuracy by training on future data (Section~\ref{sec:eval-purgedcv}).

  \item \textbf{The Deflated Sharpe Ratio} corrects backtested Sharpe ratios for the number of strategies tested. Every experiment counts as a trial \citep{Bailey2014DSR} (Section~\ref{sec:eval-dsr}).

  \item \textbf{Log every experiment.} You cannot compute an honest DSR without knowing $N$ (Section~\ref{sec:eval-dsr}).

  \item \textbf{Statistical significance is necessary but not sufficient.} Pair DM tests with economic significance: does the improvement survive transaction costs? (Section~\ref{sec:eval-pitfalls}).

  \item \textbf{Lookahead bias has four sources} in volatility pipelines: realized measures, microstructure features, options surface, and cross-asset data. Each requires a specific prevention rule; build these as hard constraints in your pipeline, not as documentation (Section~\ref{sec:eval-lookahead-taxonomy}).

  \item \textbf{Follow the full workflow} (Figure~\ref{fig:eval-workflow}): reserve holdout, log experiments, purged CV, QLIKE, MZ, DM, MCS, DSR.
\end{itemize}

\subsection{Key Results Recap}

\begin{center}
\begin{tabular}{p{3.2cm} p{4.5cm} p{5.5cm}}
\toprule
\textbf{Tool} & \textbf{What It Does} & \textbf{Key Reference} \\
\midrule
QLIKE loss & Primary loss function; proxy-robust, outlier-resistant & \citet{Patton2011} \\
MSE loss & Secondary loss; proxy-robust but outlier-sensitive & \citet{Patton2011} \\
Mincer--Zarnowitz & Diagnoses forecast bias and inefficiency & \citet{MincerZarnowitz1969} \\
Diebold--Mariano & Pairwise statistical test of loss difference & \citet{DieboldMariano1995} \\
Model Confidence Set & Multi-model comparison; returns surviving set & \citet{HansenLundeNason2011} \\
Purged K-fold CV & Time-series CV without look-ahead & \citet{LopezdePrado2018} \\
Deflated Sharpe Ratio & Corrects Sharpe for multiple testing & \citet{Bailey2014DSR} \\
Haircut Sharpe & Alternative multiple-testing correction & \citet{HarveyLiu2015} \\
\bottomrule
\end{tabular}
\end{center}
